\documentclass[10pt,a4paper]{amsart}
\usepackage[margin=19mm]{geometry}
\usepackage{amsmath,amssymb,mathtools}
\usepackage[scr=rsfs,cal=euler]{mathalfa}
\usepackage{xfrac}
\usepackage[unicode,hidelinks,bookmarks=false]{hyperref}
\newcommand{\ie}{i.e.\ }
\newcommand{\cf}{cf.\ }
\newcommand{\resp}{respectively\ }
\newcommand{\Subsection}{\subsection}
\providecommand{\nobreakdash}{\nobreak\hbox{-}}
\setlength{\emergencystretch}{2em}
\multlinegap0pt
\usepackage{bm}
\usepackage{bbm}
\usepackage{dsfont}
\usepackage{xspace}
\usepackage{cancel}
\usepackage{mathtools}
\usepackage{amsthm}
\makeatletter
\@ifundefined{theorem}{\newtheorem{theorem}{Theorem}}{}
\makeatother
\newcommand{\mfn}{\mathbf n}
\title[Generalized complex symmetric composition operators with app]{Generalized complex symmetric composition operators with applications}
\author{Vasudevarao Allu, Satyajit Sahoo}
\date{Source: arXiv:2502.20875v2 (CC BY 4.0)}
\begin{document}
\maketitle

\noindent\emph{Source and licence.} Generalized complex symmetric composition operators with applications. Version arXiv:2502.20875v2 (CC BY 4.0), distributed under the Creative Commons Attribution 4.0 International licence, as recorded in the arXiv licence metadata for this exact version and shown on the abstract page https://arxiv.org/abs/2502.20875v2. Full rights notice: licenses/arxiv-2502.20875-CC-BY-4.0.txt.\par\smallskip

\noindent\textit{Editorial scope.} Counted unit: the Theorem whose source label is \texttt{Thm\_1} in the article (source0.tex lines 251--262, source label \texttt{Thm\_1}), delivered together with the article's own complete original proof of it (source0.tex lines 263--391, one \texttt{proof} environment of 129 lines). No mathematics is written, repaired or re-derived: statement and proof are the article's own text line for line. Required context retained uncounted: none. Every definition, display and label of those lines is retained unchanged. Omitted from this package and not redistributed: 1--250, 392--1690. Nothing inside a retained excerpt is omitted or altered: every retained body is delivered exactly as the article's lines are written, so no omission receipt of any kind accompanies this package. Citations: the retained excerpts carry no citation command at all, so no bibliography entry of the article is transcribed and no reference is redistributed. Reference adaptation: none was needed and none was made; every reference inside the delivered unit resolves inside it, so no unresolved reference or citation is produced. Printed slips kept literal: no printed word, symbol, hypothesis or number of the counted statement or of its proof was altered. Layout and class: the article's own class is \texttt{amsart} with packages including amsmath, amsthm, amscd, amsfonts, amssymb, graphicx, color, hyperref, caption, subcaption, adjustbox, wrapfig, float, tikz; the wrapper is the lane's pinned \texttt{amsart} editorial wrapper at 10pt on A4, which restates the theorem-like environments the retained text needs and adds the article's own macro definitions that the retained text needs (\texttt{\textbackslash mfn}), with extra wrapper packages (bm, bbm, dsfont, xspace, cancel, mathtools). The delivered numbering is the unit's own. Assets: no retained span contains a figure environment, a graphics-inclusion command, a PDF-page inclusion, a table or a TikZ picture; no image file is copied into this package, the package declares no asset dependency and no image file is redistributed with it. Licence: version arXiv:2502.20875v2 of this article is distributed under the Creative Commons Attribution 4.0 International licence, as recorded in the arXiv licence metadata for this exact version and shown on the abstract page https://arxiv.org/abs/2502.20875v2. A full rights notice is delivered at licenses/arxiv-2502.20875-CC-BY-4.0.txt. Characters: every retained excerpt is pure ASCII, so the delivered engine, XeLaTeX with its default Latin Modern text font, reports no missing character.
\begin{theorem}\label{Thm_1}
Let $\gamma\in \mathbb{N}^+$ and $\mfn$ be a non-negative integer. Let $\psi: \mathbb{D}^d\longrightarrow \mathbb{C}$ be a non-zero analytic function and $\varphi: \mathbb{D}^d\longrightarrow \mathbb{D}^d$ with $\varphi(z_1,  \dots, z_d)=(\varphi_1(z_1), \dots, \varphi_d(z_d))$ be an analytic self map of $\mathbb{D}^d$ such that $D_{\mfn,\psi,\varphi}$ is bounded on $\mathcal{H}_\gamma(\mathbb{D}^d)$.
    Then $D_{\mfn,\psi,\varphi}$ is complex symmetric with respect to the conjugation $\mathcal{J}$ if, and only if,
\begin{gather*}
    \psi(z_1, \dots, z_d)=a\prod_{j=1}^dz_j^{n_j}\big(1-z_j\varphi_j(0)\big)^{-\gamma-n_j},
\end{gather*}
and 
\begin{gather*}
        \varphi_t(z_t)=\varphi_t(0)+\big(1-z_t\varphi_t(0)\big)^{-1}z_t\varphi_t^{\prime}(0),~ \mbox{for}~ t=1, \dots, d,
\end{gather*}
where $a=\partial^n\psi(0, \dots, 0)$.
\end{theorem}

\begin{proof}
        Suppose that $D_{\mfn,\psi,\varphi}$ is complex symmetric with respect to the conjugation $\mathcal{J}$, then we have
    \begin{align}\label{Them_Eqn_1}
         D_{\mfn,\psi, \varphi}\mathcal{J}K_w(z)= \mathcal{J}D_{\mfn,\psi, \varphi}^*K_w(z)
    \end{align}
    for all $z=(z_1, \dots, z_d)$, $w=(w_1, \dots, w_d)\in \mathbb{D}^d$.
    The left hand side of \eqref{Them_Eqn_1} gives us 
    \begin{align}\label{Them_Eqn_2}
         D_{\mfn,\psi, \varphi}\mathcal{J}K_w(z)&=\psi(z)\partial^n\mathcal{J}K_w(\varphi(z))\nonumber\\
         &=\psi(z_1, \dots, z_d)\partial^n\mathcal{J}K_{(w_1, \dots, w_d)}(\varphi_1(z_1), \dots, \varphi_d(z_d))\nonumber\\
         &=\psi(z_1, \dots, z_d)\partial^n\overline{K_{(w_1, \dots, w_d)}(\overline{\varphi_1(z_1)}, \dots, \overline{\varphi_d(z_d))}}\nonumber\\
          &=\psi(z_1, \dots, z_d)\partial^n\overline{K_{w_1}(\overline{\varphi_1(z_1)})},\dots, \overline{K_{w_d}(\overline{\varphi_d(z_d)})}\nonumber\\
           &=\psi(z_1, \dots, z_d)\partial^n\prod_{j=1}^d(1-w_j\varphi_j(z_j))^{-\gamma}\nonumber\\
            &=\psi(z_1, \dots, z_d)\prod_{j=1}^d(\gamma+n_j-1)!w_j^{n_j}(1-w_j\varphi_j(z_j))^{-\gamma-n_j}.
    \end{align}
    On the other hand, we have 
    \begin{align}\label{Them_Eqn_3}
    	\mathcal{J} D_{\mfn,\psi, \varphi}^*K_w(z)&=\mathcal{J}\overline{\psi(w)}K_{\varphi(w)}^{[\mfn]}(z)\nonumber\\
    	&=\mathcal{J}\overline{\psi(w)}\prod_{j=1}^d(\gamma+n_j-1)!z_j^{n_j}(1-z_j\overline{\varphi_j(w_j)})^{-\gamma-n_j}\nonumber\\
    	&=\psi(w_1, \dots, w_d)\prod_{j=1}^d(\gamma+n_j-1)!z_j^{n_j}(1-z_j\varphi_j(w_j))^{-\gamma-n_j}.
    \end{align}
    Hence from \eqref{Them_Eqn_1}, we have 
    \begin{align}\label{Them_Eqn_4}
    \psi(z_1, \dots, z_d)&\prod_{j=1}^d(\gamma+n_j-1)!w_j^{n_j}(1-w_j\varphi_j(z_j))^{-\gamma-n_j}\nonumber\\
    &=\psi(w_1, \dots, w_d)\prod_{j=1}^d(\gamma+n_j-1)!z_j^{n_j}(1-z_j\varphi_j(w_j))^{-\gamma-n_j}.
    \end{align}
    Letting $z=0, i.e., z_1=0,\dots,z_d=0$ in \eqref{Them_Eqn_4}, we get $\psi(0)=0$, that means $\psi(0, z_2, \dots,z_d)=0, \psi(z_1, 0 \dots,z_d)=0,\dots,\psi(z_1, z_2 \dots,z_{d-1}, 0)=0 $,  $\psi(0, \dots, 0)=0$ and so we can write $\psi(z)=g(z_1, \dots, z_d)\prod_{j=1}^dz_j^{b_j}$ and $g$ is analytic with $g(0, z_2, \dots, z_d)\neq 0, g(z_1, 0, \dots, z_d)\neq 0,\dots, g(z_1, \dots, z_{d-1},0)\neq 0$, $g(0,  \dots, 0)\neq 0$ and $b_j\in \mathbb{N}, j\in \{ 
    1, \dots, d\}$.
    Then, we get from \eqref{Them_Eqn_4} that
    \begin{align}\label{Them_Eqn_5}
   g(z_1, \dots, z_d)&\prod_{j=1}^d(\gamma+n_j-1)!w_j^{n_j}z_j^{b_j}(1-w_j\varphi_j(z_j))^{-\gamma-n_j}\nonumber\\
    &=g(w_1, \dots, w_d)\prod_{j=1}^d(\gamma+n_j-1)!w_j^{b_j}z_j^{n_j}(1-z_j\varphi_j(w_j))^{-\gamma-n_j}.
    \end{align}
    We show that $n_j=b_j$ for $j\in \{ 
    1, \dots, d\}$.\\ \\
    \underline{\bf Case-I:} If $n_j>b_j$, then it follows from \eqref{Them_Eqn_5} that 
     \begin{align}\label{Them_Eqn_6}
    g(z_1, \dots, z_d)&\prod_{j=1}^dw_j^{n_j-b_j}(1-w_j\varphi_j(z_j))^{-\gamma-n_j}\nonumber\\
   & =g(w_1, \dots, w_d)\prod_{j=1}^dz_j^{n_j-b_j}(1-z_j\varphi_j(w_j))^{-\gamma-n_j},
    \end{align}
    which is impossible because by letting $z_1=0,\dots,z_d=0$ in \eqref{Them_Eqn_6}, we get    $g(0, z_2, \dots, z_d)= 0, g(z_1, 0, \dots, z_d)= 0,\dots, g(z_1, \dots, z_{d-1},0)= 0$ and $g(0,  \dots, 0)= 0$ which is a contradiction. This contradiction is due to the fact that  $g(0, z_2, \dots, z_d)\neq 0, g(z_1, 0, \dots, z_d)\neq 0,\dots, g(z_1, \dots, z_{d-1},0)\neq 0$ and $g(0,  \dots, 0)\neq 0$. \\ \\
    \underline{\bf Case-II:} If $n_j<b_j$, then \eqref{Them_Eqn_5} becomes the following
     \begin{align}\label{Them_Eqn_7}
    g(z_1, \dots, z_d)&\prod_{j=1}^dz_j^{b_j-n_j}(1-w_j\varphi_j(z_j))^{-\gamma-n_j}\nonumber\\
    &=g(w_1, \dots, w_d)\prod_{j=1}^dw_j^{b_j-n_j}(1-z_j\varphi_j(w_j))^{-\gamma-n_j},
    \end{align}
    which is not possible because by letting $w_1=0,\dots,w_d=0$ in \eqref{Them_Eqn_7}, we get    $g(z,  \dots, z_d)\equiv 0$ on $\mathbb{D}^d$ which is a contradiction. Thus, we must have $n_j=b_j$ for all  $j\in \{ 
    1, \dots, d\}$ and hence \eqref{Them_Eqn_5} reduces to the following   
     \begin{align}\label{Them_Eqn_8}
    g(z_1, \dots, z_d)\prod_{j=1}^d(1-w_j\varphi_j(z_j))^{-\gamma-n_j}
    =g(w_1, \dots, w_d)\prod_{j=1}^d(1-z_j\varphi_j(w_j))^{-\gamma-n_j}.
    \end{align}
    Letting $w_j=0$ for all $j\in \{ 
    1, \dots, d\}$ in \eqref{Them_Eqn_8}, we obtain
    \begin{gather}\label{Them_Eqn_9}
   g(z_1, \dots, z_d)=g(0,  \dots, 0)\prod_{j=1}^d\big(1-z_j\varphi_j(0)\big)^{-\gamma-n_j}.
    \end{gather}
    It follows that 
    \begin{align}\label{Them_Eqn_10}
    \psi(z_1, \dots, z_d)&=g(z_1, \dots, z_d)\prod_{j=1}^dz_j^{n_j}\nonumber\\
    &=g(0,  \dots, 0)\prod_{j=1}^dz_j^{n_j}\big(1-z_j\varphi_j(0)\big)^{-\gamma-n_j}.
    \end{align}
  Since $g(0,  \dots, 0)\neq 0$, it follows from \eqref{Them_Eqn_4} and \eqref{Them_Eqn_10} that 
  \begin{align}\label{Them_Eqn_11}
  	\prod_{j=1}^d\big(1-w_j\varphi_j(0)\big)^{\gamma+n_j}\times	&\prod_{j=1}^d\big(1-z_j\varphi_j(w_j)\big)^{\gamma+n_j}\nonumber\\
  	&=\prod_{j=1}^d\big(1-z_j\varphi_j(0)\big)^{\gamma+n_j}\times	\prod_{j=1}^d\big(1-w_j\varphi_j(z_j)\big)^{\gamma+n_j}.
  	  \end{align}
  	  Through differentiating both sides with respect to the variable $z_t$, $(t\in\{1, \dots, d\})$ in \eqref{Them_Eqn_11} we have
  	  \begin{align}\label{Them_Eqn_12}
  	& (\gamma+n_t) \prod_{j=1}^d\big(1-w_j\varphi_j(0)\big)^{\gamma+n_j}\times	\prod_{j\neq t}^d\big(1-z_j\varphi_j(w_j)\big)^{\gamma+n_j}\big(1-z_t\varphi_t(w_t)\big)^{\gamma+n_t-1}(-\varphi_t(w_t))\nonumber\\
  	 &= (\gamma+n_t)\prod_{j\neq t}^d\big(1-z_j\varphi_j(0)\big)^{\gamma+n_j}\big(1-z_t\varphi_t(0)\big)^{\gamma+n_t-1}(-\varphi_t(0))\times	\prod_{j=1}^d\big(1-w_j\varphi_j(z_j)\big)^{\gamma+n_j}\nonumber\\
  	 &\hspace{2cm}+(\gamma+n_t)\prod_{j=1}^d\big(1-z_j\varphi_j(0)\big)^{\gamma+n_j}\prod_{j\neq t}^d\big(1-w_j\varphi_j(z_j)\big)^{\gamma+n_j}\nonumber\\
     &\hspace{7cm}\times	\big(1-w_t\varphi_t(z_t)\big)^{\gamma+n_t-1}\times(-w_t\varphi_t^{\prime}(z_t)).
  	  \end{align}
  	  Evaluating at $z_j=0$ for all $j=1, \dots, d$ and $w_j=0$ for $j\neq t$ in \eqref{Them_Eqn_12}, we obtain 
  	  \begin{align}\label{Them_Eqn_13}
  	  	\big(1-w_t\varphi_t(0)\big)^{\gamma+n_t}&(\varphi_t(w_t))\nonumber\\
  	  	&=	\big(1-w_t\varphi_t(0)\big)^{\gamma+n_t}(\varphi_t(0))+\big(1-w_t\varphi_t(0)\big)^{\gamma+n_t-1}(w_t\varphi_t^{\prime}(0)),
  	  \end{align} 
  	  which implies that 
  	  \begin{gather*}
  	  \varphi_t(w_t)=\varphi_t(0)+\big(1-w_t\varphi_t(0)\big)^{-1}w_t\varphi_t^{\prime}(0),~ \mbox{for}~ t=1, \dots, d.
  	  \end{gather*}
  	  Thus,
  	  \begin{gather*}
  	  \psi(z_1, \dots, z_d)=a\prod_{j=1}^dz_j^{n_j}\big(1-z_j\varphi_j(0)\big)^{-\gamma-n_j},
  	  \end{gather*}
  	  and 
  	  \begin{gather*}
  	  \varphi_t(z_t)=\varphi_t(0)+\big(1-z_t\varphi_t(0)\big)^{-1}z_t\varphi_t^{\prime}(0),~ \mbox{for}~ t=1, \dots, d,
  	  \end{gather*}
  	  where $$a=\partial^n\psi(0, \dots, 0)=\frac{\partial^{n_1+\cdots+n_d}}{\partial z_1^{n_1},\dots, \partial z_d^{n_d}}\psi(0, \dots, 0).$$
  	  
  	  Conversly assume that 
  	  \begin{gather*}
  	  \psi(z_1, \dots, z_d)=a\prod_{j=1}^dz_j^{n_j}\big(1-z_j\varphi_j(0)\big)^{-\gamma-n_j},
  	  \end{gather*}
  	  and 
  	  \begin{gather*}
  	  \varphi_t(z_t)=\varphi_t(0)+\big(1-z_t\varphi_t(0)\big)^{-1}z_t\varphi_t^{\prime}(0),~ \mbox{for}~ t=1, \dots, d,
  	  \end{gather*}
  	  where $a=\partial^n\psi(0, \dots, 0)$.
  	  From \eqref{Them_Eqn_2} and \eqref{Them_Eqn_3}, we have  
  	  
  	   \begin{align}\label{Them_Eqn_14}
  	   D_{\mfn,\psi, \varphi}&\mathcal{J}K_w(z)
  	 = \psi(z_1, \dots, z_d)\prod_{j=1}^d(\gamma+n_j-1)!w_j^{n_j}(1-w_j\varphi_j(z_j))^{-\gamma-n_j}\nonumber\\
  	  =&a\prod_{j=1}^dz_j^{n_j}\big(1-z_j\varphi_j(0)\big)^{-\gamma-n_j}\times\prod_{j=1}^d(\gamma+n_j-1)!w_j^{n_j}\nonumber\\
  	 ~~~~~~~~~~~~~~~~~~~~&\left(1-w_j\left(\varphi_j(0)+\big(1-z_j\varphi_j(0)\big)^{-1}z_j\varphi_j^{\prime}(0)\right)\right)^{-\gamma-n_j}\nonumber\\
  	  =&a\prod_{j=1}^d(\gamma+n_j-1)!z_j^{n_j}w_j^{n_j}\left(1-z_j\varphi_j(0)-w_j\left((1-z_j\varphi_j(0))\varphi_j(0)+z_j\varphi_j^{\prime}(0)\right) \right)^{-\gamma-n_j}.
  	  \end{align} 
  	  and 
  	   \begin{align}\label{Them_Eqn_15}
  	  &  \mathcal{J}D_{\mfn,\psi, \varphi}^*K_w(z)
  	  = \psi(w_1, \dots, w_d)\prod_{j=1}^d(\gamma+n_j-1)!z_j^{n_j}(1-z_j\varphi_j(w_j))^{-\gamma-n_j}\nonumber\\
  	  =&a\prod_{j=1}^dw_j^{n_j}\big(1-w_j\varphi_j(0)\big)^{-\gamma-n_j}\times\prod_{j=1}^d(\gamma+n_j-1)!z_j^{n_j}\nonumber\\
  	  &\left(1-z_j\left(\varphi_j(0)+\big(1-w_j\varphi_j(0)\big)^{-1}w_j\varphi_j^{\prime}(0)\right)\right)^{-\gamma-n_j}\nonumber\\
  	  =&a\prod_{j=1}^d(\gamma+n_j-1)!z_j^{n_j}w_j^{n_j}\left(1-w_j\varphi_j(0)-z_j\left((1-w_j\varphi_j(0))\varphi_j(0)+w_j\varphi_j^{\prime}(0)\right) \right)^{-\gamma-n_j},
  	  \end{align} 
  	   for all $z=(z_1, \dots, z_d)$, $w=(w_1, \dots, w_d)\in \mathbb{D}^d$.
  	   Since 
  	   \begin{align}\label{Them_Eqn_16}
  	   &1-z_j\varphi_j(0)-w_j\left((1-z_j\varphi_j(0))\varphi_j(0)+z_j\varphi_j^{\prime}(0)\right)\nonumber\\
  	   &=1-z_j\varphi_j(0)-w_j\varphi_j(0)-w_jz_j\varphi_j^2(0)+w_jz_j\varphi_j^{\prime}(0)\nonumber\\
  	   &=1-w_j\varphi_j(0)-z_j\left((1-w_j\varphi_j(0))\varphi_j(0)+w_j\varphi_j^{\prime}(0)\right),
  	   \end{align}
  	   it follows from \eqref{Them_Eqn_14} and \eqref{Them_Eqn_15} that 
  	   $D_{\mfn,\psi, \varphi}\mathcal{J}K_w(z)=\mathcal{J}D_{\mfn,\psi, \varphi}^*K_w(z)$. Hence $D_{\mfn,\psi,\varphi}$ is complex symmetric with respect to the conjugation $\mathcal{J}$. This completes the proof.
\end{proof}

\end{document}
